\documentclass[11pt,a4paper]{article}

\usepackage[utf8]{inputenc}
\usepackage[T1]{fontenc}
\usepackage[margin=2.5cm]{geometry}
\usepackage{hyperref}
\usepackage{xcolor}
\usepackage{listings}
\usepackage{enumitem}
\usepackage{booktabs}
\usepackage{longtable}
\usepackage{amsmath,amssymb}
\usepackage{fancyhdr}
\usepackage{titlesec}

\definecolor{isls-blue}{HTML}{1a5276}
\definecolor{isls-orange}{HTML}{e67e22}
\definecolor{isls-green}{HTML}{27ae60}
\definecolor{isls-red}{HTML}{c0392b}
\definecolor{isls-gray}{HTML}{7f8c8d}
\definecolor{code-bg}{HTML}{f7f9fb}

\lstdefinestyle{rust}{
  language=C,
  basicstyle=\ttfamily\small,
  keywordstyle=\color{isls-blue}\bfseries,
  commentstyle=\color{isls-gray}\itshape,
  stringstyle=\color{isls-green},
  backgroundcolor=\color{code-bg},
  frame=single,
  rulecolor=\color{isls-gray!30},
  breaklines=true,
  showstringspaces=false,
  tabsize=4,
  morekeywords={fn,let,mut,pub,struct,impl,enum,use,mod,self,Self,
    crate,super,where,trait,for,in,if,else,match,return,async,await,
    Ok,Err,Some,None,Vec,String,HashMap,HashSet,PathBuf,Result,Option,
    bool,true,false,usize},
}
\lstdefinestyle{bash}{
  language=bash,
  basicstyle=\ttfamily\small,
  backgroundcolor=\color{code-bg},
  frame=single,
  rulecolor=\color{isls-gray!30},
  breaklines=true,
  showstringspaces=false,
}
\lstdefinestyle{toml}{
  basicstyle=\ttfamily\small,
  backgroundcolor=\color{code-bg},
  frame=single,
  rulecolor=\color{isls-gray!30},
  breaklines=true,
  showstringspaces=false,
  commentstyle=\color{isls-gray}\itshape,
  morecomment=[l]{\#},
}

\pagestyle{fancy}
\fancyhf{}
\fancyhead[L]{\small\textcolor{isls-gray}{ISLS D6 Spec --- M\"obius-Umst\"ulpung}}
\fancyhead[R]{\small\textcolor{isls-gray}{v2.0.0 --- \today}}
\fancyfoot[C]{\thepage}

\titleformat{\section}
  {\Large\bfseries\color{isls-blue}}{\thesection}{1em}{}
\titleformat{\subsection}
  {\large\bfseries\color{isls-blue!80}}{\thesubsection}{1em}{}
\titleformat{\subsubsection}
  {\normalsize\bfseries\color{isls-blue!60}}{\thesubsubsection}{1em}{}

\hypersetup{
  colorlinks=true,
  linkcolor=isls-blue,
  urlcolor=isls-orange,
  citecolor=isls-green,
}

% ══════════════════════════════════════════════════════════════════════════════
\begin{document}

\begin{titlepage}
\centering
\vspace*{2.5cm}

{\Huge\bfseries\color{isls-blue} ISLS D6 Specification}\\[0.5cm]
{\LARGE\color{isls-orange} M\"obius-Umst\"ulpung:\\Norm-getriebener HDAG \& Infrastruktur als emergente Eigenschaft}\\[2cm]

{\large
\begin{tabular}{rl}
\textbf{System:}    & ISLS --- Intelligent Semantic Ledger Substrate \\
\textbf{Version:}   & v4.0 (D6) \\
\textbf{Author:}    & Sebastian Klemm \\
\textbf{Date:}      & \today \\
\textbf{Status:}    & Specification \\
\textbf{Depends:}   & D5 (complete), all 10 crates \\
\textbf{Branch:}    & \texttt{claude/implement-isls-d6-spec} \\
\end{tabular}
}

\vfill

{\small\color{isls-gray}
Mithras Substructure University\\
Repository: \texttt{https://github.com/LashSesh/graphity}\\
License: MIT
}
\end{titlepage}

\tableofcontents
\newpage


% ══════════════════════════════════════════════════════════════════════════════
\section{Executive Summary}
% ══════════════════════════════════════════════════════════════════════════════

ISLS generates applications. Until now, it generates exactly one kind: 
Actix-Web/PostgreSQL/Docker CRUD web apps. This is hardcoded in 
\texttt{CodegenHdag::build()}, which constructs a fixed file tree
(48 structural + 7 LLM nodes) regardless of what the user describes.

D6 removes this constraint. The HDAG becomes \textbf{norm-driven}: 
the file tree emerges from the composition of activated norms, not
from a hardcoded function. A web-app description activates 
web-infrastructure norms $\to$ web-app HDAG. A CLI-tool description
activates CLI-infrastructure norms $\to$ CLI-tool HDAG. A library
description activates library-infrastructure norms $\to$ library HDAG.
No categories defined. No archetypes. The infrastructure is an
\emph{emergent property} of the description matched against the norm 
catalog.

The M\"obius proof: ISLS generates a management dashboard for itself
(ISLS Studio), demonstrating that the same pipeline produces different
infrastructure configurations from different descriptions.

\subsection{The Architectural Problem}

\texttt{hdag.rs::CodegenHdag::build()} contains $\sim$200 lines that
hardcode:

\begin{itemize}[nosep]
\item \texttt{backend/Cargo.toml} with actix-web, sqlx, jsonwebtoken
\item \texttt{docker-compose.yml} with PostgreSQL, Nginx
\item \texttt{backend/src/models/}, \texttt{database/}, \texttt{services/}, 
      \texttt{api/} directories
\item \texttt{frontend/} with index.html, style.css, client.js, 
      entity pages
\item \texttt{backend/Dockerfile} with \texttt{rust:1.88-slim}
\item Fixed layer ordering: 0 (static) $\to$ 1 (foundation) $\to$ 
      2 (auth) $\to$ 3 (models) $\to$ 4 (database) $\to$ 5 (services) 
      $\to$ 6 (api) $\to$ 7 (main) $\to$ 8 (frontend) $\to$ 9 (tests)
\end{itemize}

Every node, every edge, every \texttt{ProvidedSymbol} assumes this 
specific shape. Generating a CLI tool, a library, or a desktop app 
is impossible without rewriting the entire function.

\subsection{The Solution: InfraBlueprint}

An \textbf{InfraBlueprint} is a pure-data struct that describes which 
file categories, layers, and infrastructure components a project has.
It is derived from the activated norms at generation time. 
\texttt{CodegenHdag::build()} takes \texttt{(AppSpec, InfraBlueprint)} 
instead of just \texttt{AppSpec} and constructs nodes/edges according 
to the blueprint.

The current web-app structure becomes the output of composing a specific
set of infrastructure norms. Not a hardcoded default. Just the first 
norm family that was built.

\subsection{What D6 Achieves}

\begin{enumerate}[nosep]
\item \textbf{InfraBlueprint} data structure and derivation from norms.
\item \textbf{Infrastructure Norms:} 7 new builtin norms defining 
      web-server, database, auth, frontend, docker, CLI, and library 
      infrastructure.
\item \textbf{Refactored HDAG:} \texttt{CodegenHdag::build()} reads 
      blueprint, not hardcoded assumptions.
\item \textbf{Refactored Structural Generators:} Dispatched by 
      blueprint capabilities, not by path-matching.
\item \textbf{ISLS Studio TOML + forge-self:} Self-description that 
      activates web-infrastructure norms through description matching.
\item \textbf{Backward Compatibility:} All D1--D5 apps (warehouse, 
      hotel, clinic, gym, school, spa) still generate identically.
\end{enumerate}


% ══════════════════════════════════════════════════════════════════════════════
\section{InfraBlueprint}
\label{sec:blueprint}
% ══════════════════════════════════════════════════════════════════════════════

\subsection{Data Structure}

\begin{lstlisting}[style=rust]
/// Infrastructure blueprint derived from activated norms.
/// Describes WHAT the generated project contains — not HOW 
/// to generate it (that's the structural generators' job).
#[derive(Clone, Debug, Default)]
pub struct InfraBlueprint {
    // ── Output shape ─────────────────────────────────
    /// Project produces a runnable binary.
    pub has_binary: bool,
    /// Project produces a library crate (pub API).
    pub has_library: bool,
    
    // ── Server ───────────────────────────────────────
    /// Project includes an HTTP server.
    pub has_http_server: bool,
    /// Server framework crate name (e.g. "actix-web").
    pub server_framework: Option<String>,
    
    // ── Persistence ──────────────────────────────────
    /// Project uses a database.
    pub has_database: bool,
    /// Database type (e.g. "postgresql", "sqlite").
    pub database_type: Option<String>,
    /// Database driver crate (e.g. "sqlx").
    pub database_driver: Option<String>,
    
    // ── Auth ─────────────────────────────────────────
    /// Project includes authentication.
    pub has_auth: bool,
    /// Auth mechanism (e.g. "jwt").
    pub auth_type: Option<String>,
    
    // ── Frontend ─────────────────────────────────────
    /// Project includes a web frontend.
    pub has_frontend: bool,
    /// Frontend type (e.g. "vanilla-js", "react").
    pub frontend_type: Option<String>,
    
    // ── Deployment ───────────────────────────────────
    /// Project includes Docker files.
    pub has_docker: bool,
    
    // ── CLI ──────────────────────────────────────────
    /// Project includes CLI argument parsing.
    pub has_cli: bool,
    /// CLI framework (e.g. "clap").
    pub cli_framework: Option<String>,
    
    // ── Crate dependencies ───────────────────────────
    /// Additional Cargo.toml dependencies from norms.
    pub extra_dependencies: Vec<CrateDep>,
}

#[derive(Clone, Debug)]
pub struct CrateDep {
    pub name: String,
    pub version: String,
    pub features: Vec<String>,
}
\end{lstlisting}

\subsection{Key Principle: Booleans, Not Enums}

The blueprint uses boolean flags (\texttt{has\_http\_server}, 
\texttt{has\_database}, \texttt{has\_cli}), not an enum like 
\texttt{ProjectType::WebApp}. This is deliberate:

\begin{itemize}[nosep]
\item A project can have BOTH an HTTP server AND a CLI 
      (e.g. a web server with CLI management commands).
\item A project can have a database WITHOUT an HTTP server 
      (e.g. a batch-processing CLI tool with PostgreSQL).
\item A project can have NEITHER server nor CLI --- just a library.
\item The blueprint describes \emph{capabilities}, not \emph{categories}.
\end{itemize}

No combination is invalid. The HDAG builder creates nodes for whatever
capabilities the blueprint declares. If a capability is absent, the
corresponding nodes are simply not created.

\subsection{Derivation from Norms}

\begin{lstlisting}[style=rust]
/// Derive an InfraBlueprint from activated infrastructure norms.
pub fn derive_blueprint(
    activated: &[ActivatedNorm],
) -> InfraBlueprint {
    let mut bp = InfraBlueprint::default();
    
    for norm in activated {
        match norm.norm.id.as_str() {
            "ISLS-NORM-INFRA-WEB" => {
                bp.has_binary = true;
                bp.has_http_server = true;
                bp.server_framework = Some("actix-web".into());
            }
            "ISLS-NORM-INFRA-DB" => {
                bp.has_database = true;
                bp.database_type = Some("postgresql".into());
                bp.database_driver = Some("sqlx".into());
            }
            "ISLS-NORM-INFRA-AUTH" => {
                bp.has_auth = true;
                bp.auth_type = Some("jwt".into());
            }
            "ISLS-NORM-INFRA-FRONTEND" => {
                bp.has_frontend = true;
                bp.frontend_type = Some("vanilla-js".into());
            }
            "ISLS-NORM-INFRA-DOCKER" => {
                bp.has_docker = true;
            }
            "ISLS-NORM-INFRA-CLI" => {
                bp.has_binary = true;
                bp.has_cli = true;
                bp.cli_framework = Some("clap".into());
            }
            "ISLS-NORM-INFRA-LIB" => {
                bp.has_library = true;
            }
            _ => {} // Entity/feature norms don't affect blueprint
        }
    }
    
    // Default: if nothing activated, assume web app
    // (backward compat with D1-D5 TOML files)
    if !bp.has_binary && !bp.has_library {
        bp.has_binary = true;
        bp.has_http_server = true;
        bp.server_framework = Some("actix-web".into());
        bp.has_database = true;
        bp.database_type = Some("postgresql".into());
        bp.database_driver = Some("sqlx".into());
        bp.has_auth = true;
        bp.auth_type = Some("jwt".into());
        bp.has_frontend = true;
        bp.frontend_type = Some("vanilla-js".into());
        bp.has_docker = true;
    }
    
    bp
}
\end{lstlisting}

The default fallback ensures all existing D1--D5 apps generate 
identically. The system only deviates from the web-app shape when 
infrastructure norms explicitly configure it differently.


% ══════════════════════════════════════════════════════════════════════════════
\section{Infrastructure Norms}
\label{sec:infra-norms}
% ══════════════════════════════════════════════════════════════════════════════

Seven new builtin norms are added to \texttt{catalog.rs}. They have 
\texttt{NormLevel::Atom} (infrastructure primitives, below Molecule).

\begin{longtable}{llp{6cm}}
\toprule
\textbf{ID} & \textbf{Name} & \textbf{Triggers / Keywords} \\
\midrule
\texttt{ISLS-NORM-INFRA-WEB} & Web-Server 
  & web, server, api, rest, http, endpoint, backend, app, 
    application, service, microservice \\
\texttt{ISLS-NORM-INFRA-DB} & Database 
  & database, db, postgres, sql, storage, persist, table, 
    query, migration, data \\
\texttt{ISLS-NORM-INFRA-AUTH} & Authentication 
  & auth, login, register, jwt, token, user, password, 
    session, security \\
\texttt{ISLS-NORM-INFRA-FRONTEND} & Frontend 
  & frontend, ui, page, dashboard, interface, browser, web ui, 
    html, css, javascript \\
\texttt{ISLS-NORM-INFRA-DOCKER} & Docker 
  & docker, container, deploy, compose, production, 
    deployment \\
\texttt{ISLS-NORM-INFRA-CLI} & CLI-Tool 
  & cli, command, terminal, tool, script, batch, utility, 
    convert, process, parse \\
\texttt{ISLS-NORM-INFRA-LIB} & Library 
  & library, crate, lib, sdk, api (without ``server''/``web''),
    module, package \\
\bottomrule
\end{longtable}

\paragraph{Activation Logic.} Infrastructure norms are matched against 
the user's description by the same \texttt{match\_description()} 
function as entity norms. The description ``Restaurant app with table 
reservations'' triggers WEB (``app''), DB (``table'', ``data''), AUTH 
(implicitly via ``app'' $\to$ users), FRONTEND (``app'' implies UI), 
DOCKER (``app'' implies deployment). The description ``CLI tool for 
converting CSV files'' triggers CLI (``cli'', ``tool'', ``convert'')
and possibly DB if ``database'' appears.

\paragraph{NormLayers.} Infrastructure norms have empty 
\texttt{NormLayers} --- they carry no entity-level artifacts. Their 
purpose is to influence the \texttt{InfraBlueprint}, not to inject 
database tables or models.

\paragraph{Excludes.} ISLS-NORM-INFRA-LIB has 
\texttt{excludes: ["server", "web", "app", "deploy"]} to prevent 
activation alongside web-server norms (a library doesn't serve HTTP).


% ══════════════════════════════════════════════════════════════════════════════
\section{Refactored HDAG Builder}
\label{sec:hdag-refactor}
% ══════════════════════════════════════════════════════════════════════════════

\subsection{Current: \texttt{CodegenHdag::build(spec)}}

\begin{lstlisting}[style=rust]
// BEFORE: 200 lines of hardcoded nodes
pub fn build(spec: &AppSpec) -> Self {
    // Layer 0: ALWAYS creates Cargo.toml, Dockerfile, 
    //          docker-compose, nginx, frontend...
    // Layer 1: ALWAYS creates errors.rs, pagination.rs, 
    //          auth.rs, config.rs, pool.rs...
    // ... etc.
}
\end{lstlisting}

\subsection{Target: \texttt{CodegenHdag::build(spec, blueprint)}}

\begin{lstlisting}[style=rust]
pub fn build(spec: &AppSpec, bp: &InfraBlueprint) -> Self {
    let mut nodes = Vec::new();
    let mut edges = Vec::new();
    // helper closure: add_node(...)
    
    // ── Layer 0: Infrastructure (conditional) ─────────
    add_node("backend/Cargo.toml", Structural, 0, ...);
    add_node(".gitignore", Structural, 0, ...);
    
    if bp.has_docker {
        add_node("backend/Dockerfile", Structural, 0, ...);
        add_node("docker-compose.yml", Structural, 0, ...);
        add_node(".env.example", Structural, 0, ...);
    }
    if bp.has_frontend {
        add_node("frontend/nginx.conf", Structural, 0, ...);
        add_node("frontend/index.html", Structural, 8, ...);
        add_node("frontend/style.css", Structural, 8, ...);
        add_node("frontend/src/api/client.js", Structural, 8,...);
        for entity in non_user_entities {
            add_node(frontend_page(entity), Structural, 8,...);
        }
    }
    
    // ── Layer 1: Foundation (conditional) ─────────────
    if bp.has_http_server {
        errors_idx = add_node("backend/src/errors.rs",...);
        config_idx = add_node("backend/src/config.rs",...);
    }
    if bp.has_http_server && bp.has_database {
        pagination_idx = add_node("backend/src/pagination.rs",...);
        pool_idx = add_node("backend/src/database/pool.rs",...);
    }
    
    // ── Layer 2: Auth (conditional) ───────────────────
    if bp.has_auth {
        user_model_idx = add_node("backend/src/models/user.rs",...);
        auth_idx = add_node("backend/src/auth.rs", ...);
    }
    
    // ── Layer 3: Models (always, if entities exist) ───
    if bp.has_database || bp.has_http_server {
        for entity in non_user_entities {
            add_node(model_path(entity), Structural, 3, ...);
        }
    }
    
    // ── Layer 4: Database (conditional) ───────────────
    if bp.has_database {
        add_node("backend/migrations/001_initial.sql",...);
        for entity in all_entities {
            add_node(query_path(entity), Llm, 4, ...);
        }
    }
    
    // ── Layer 5: Services (conditional) ───────────────
    if bp.has_http_server || bp.has_cli {
        for entity in all_entities {
            add_node(service_path(entity), Structural, 5,...);
        }
    }
    
    // ── Layer 6: API routes (conditional) ─────────────
    if bp.has_http_server {
        if bp.has_auth {
            add_node("backend/src/api/auth_routes.rs",...);
        }
        for entity in non_user_entities {
            add_node(api_route_path(entity), Structural, 6,...);
        }
    }
    
    // ── Layer 7: Entry point ─────────────────────────
    add_node("backend/src/main.rs", Structural, 7, ...);
    // main.rs content adapts to blueprint:
    // - if has_http_server: actix-web HttpServer::new()
    // - if has_cli: clap App::new() + subcommands
    // - if has_library: no main.rs (lib.rs only)
    
    // ── Layer 8: Frontend (handled above) ────────────
    
    // ── Layer 9: Tests ───────────────────────────────
    add_node("backend/tests/api_tests.rs", Structural, 9,...);
    
    // ── Mod files (always) ───────────────────────────
    // mod.rs files are created for whatever directories
    // were populated by the conditional blocks above.
    add_mod_files(&nodes, &mut edges, ...);
    
    // ── Edges ────────────────────────────────────────
    // Wired based on which nodes actually exist.
    wire_edges(&nodes, &mut edges, bp, spec);
    
    CodegenHdag { nodes, edges }
}
\end{lstlisting}

\subsection{Key Change: Conditional Node Creation}

The critical difference: nodes are only created when the blueprint
declares the corresponding capability. If \texttt{has\_frontend == false},
no frontend nodes exist. If \texttt{has\_docker == false}, no Dockerfile
or docker-compose. If \texttt{has\_auth == false}, no user model, no 
auth.rs, no auth\_routes.

The edge-wiring function \texttt{wire\_edges()} must handle missing 
nodes gracefully. If \texttt{pagination\_idx} doesn't exist because 
there's no HTTP server, the edges from pagination to queries are simply 
not created. The LLM prompts for query files then don't include 
PaginationParams --- which is correct for a non-web context.

\subsection{Structural Generator Adaptation}

\texttt{structural.rs} and \texttt{static\_files.rs} generators don't 
need fundamental changes. They are dispatched by file path, and the HDAG 
only creates paths that the blueprint permits. What changes:

\begin{itemize}[nosep]
\item \texttt{generate\_cargo\_toml()} reads \texttt{InfraBlueprint} 
      to determine which dependencies to include (actix-web? clap? both?).
\item \texttt{generate\_main\_rs()} reads blueprint to determine entry 
      point shape (HTTP server? CLI parser? both?).
\item \texttt{generate\_dockerfile()} --- only called if 
      \texttt{has\_docker}.
\item \texttt{generate\_docker\_compose()} --- only called if 
      \texttt{has\_docker}.
\end{itemize}

All other generators (models, services, routes, frontend) remain 
unchanged --- they are simply not invoked when their nodes don't exist.

The \texttt{ForgePlan} struct gains an \texttt{InfraBlueprint} field:

\begin{lstlisting}[style=rust]
pub struct ForgePlan {
    pub spec: AppSpec,
    pub norm_ids: Vec<String>,
    pub blueprint: InfraBlueprint,  // NEW
}
\end{lstlisting}


% ══════════════════════════════════════════════════════════════════════════════
\section{Work Packages}
% ══════════════════════════════════════════════════════════════════════════════

% ──────────────────────────────────────────────────────────────────────────────
\subsection{W1: InfraBlueprint + Infrastructure Norms}
\label{sec:w1}
% ──────────────────────────────────────────────────────────────────────────────

\subsubsection{Goal}

Define the \texttt{InfraBlueprint} struct and the 7 infrastructure 
norms. Pure data, no behavior changes yet.

\subsubsection{New Files}

\begin{longtable}{lp{10cm}}
\toprule
\textbf{File} & \textbf{Purpose} \\
\midrule
\texttt{isls-forge-llm/src/blueprint.rs}
  & \texttt{InfraBlueprint} struct, \texttt{CrateDep} struct,
    \texttt{derive\_blueprint()} function. \\
\bottomrule
\end{longtable}

\subsubsection{Modified Files}

\begin{longtable}{lp{10cm}}
\toprule
\textbf{File} & \textbf{Change} \\
\midrule
\texttt{isls-forge-llm/src/lib.rs}
  & Add \texttt{pub mod blueprint;} Add \texttt{blueprint: InfraBlueprint}
    to \texttt{ForgePlan}. \\
\texttt{isls-norms/src/catalog.rs}
  & Add 7 infrastructure norms to \texttt{builtin\_norms()}. \\
\texttt{isls-norms/src/types.rs}
  & Add \texttt{NormLevel::Atom} variant (below Molecule). \\
\bottomrule
\end{longtable}

\subsubsection{W1 Verification}

\begin{enumerate}[nosep]
\item \texttt{cargo build --workspace} --- compiles.
\item Unit test: \texttt{derive\_blueprint()} with web-infrastructure 
      norms produces blueprint with all web flags true.
\item Unit test: empty norm list produces default web-app blueprint 
      (backward compatibility).
\item Unit test: CLI infrastructure norm produces 
      \texttt{has\_cli = true, has\_http\_server = false}.
\end{enumerate}


% ──────────────────────────────────────────────────────────────────────────────
\subsection{W2: Refactor CodegenHdag::build()}
\label{sec:w2}
% ──────────────────────────────────────────────────────────────────────────────

\subsubsection{Goal}

Change \texttt{CodegenHdag::build()} from \texttt{build(spec)} to 
\texttt{build(spec, blueprint)}. All node creation becomes conditional 
on blueprint flags.

\subsubsection{Approach}

\begin{enumerate}[nosep]
\item Copy the current \texttt{build()} to 
      \texttt{build\_legacy(spec)} (safety net for testing).
\item Implement new \texttt{build(spec, blueprint)} with conditional
      node creation as described in Section~\ref{sec:hdag-refactor}.
\item Extract edge-wiring into \texttt{wire\_edges()} helper that 
      handles absent nodes.
\item All callers updated to pass blueprint.
\end{enumerate}

\subsubsection{Critical Invariant}

When called with the default web-app blueprint, 
\texttt{build(spec, bp\_web)} MUST produce the \textbf{exact same}
nodes and edges as the old \texttt{build(spec)}. 
This is verified by a test that compares both outputs.

\subsubsection{Modified Files}

\begin{longtable}{lp{10cm}}
\toprule
\textbf{File} & \textbf{Change} \\
\midrule
\texttt{isls-forge-llm/src/hdag.rs}
  & Refactor \texttt{build()} to \texttt{build(spec, blueprint)}.
    Add \texttt{wire\_edges()} helper. Keep 
    \texttt{build\_legacy()} for comparison test. \\
\texttt{isls-forge-llm/src/staged\_closure.rs}
  & Pass \texttt{self.plan.blueprint} to \texttt{CodegenHdag::build()}.\\
\texttt{isls-forge-llm/src/forge.rs}
  & Pass blueprint through \texttt{LlmForge}. \\
\bottomrule
\end{longtable}

\subsubsection{W2 Verification}

\begin{enumerate}[nosep]
\item \texttt{cargo test --workspace} --- all 181+ tests pass.
\item Comparison test: \texttt{build(spec, default\_web\_blueprint()) 
      == build\_legacy(spec)} for node count, edge count, node paths.
\item Reduced blueprint test: \texttt{build(spec, no\_frontend\_bp())}
      produces fewer nodes (no frontend/ nodes).
\end{enumerate}


% ──────────────────────────────────────────────────────────────────────────────
\subsection{W3: Adaptive Structural Generators}
\label{sec:w3}
% ──────────────────────────────────────────────────────────────────────────────

\subsubsection{Goal}

Make \texttt{generate\_cargo\_toml()} and \texttt{generate\_main\_rs()}
blueprint-aware so they produce correct content for different 
infrastructure configurations.

\subsubsection{Changes}

\paragraph{\texttt{generate\_cargo\_toml(spec, blueprint)}}

\begin{itemize}[nosep]
\item If \texttt{has\_http\_server}: include actix-web, actix-cors.
\item If \texttt{has\_database}: include sqlx with appropriate feature 
      flags.
\item If \texttt{has\_auth}: include jsonwebtoken, bcrypt.
\item If \texttt{has\_cli}: include clap with derive feature.
\item If \texttt{has\_library}: \texttt{[lib]} section instead of / 
      in addition to \texttt{[[bin]]}.
\item Always: serde, tracing, thiserror, chrono.
\end{itemize}

\paragraph{\texttt{generate\_main\_rs(spec, blueprint)}}

\begin{itemize}[nosep]
\item If \texttt{has\_http\_server}: Actix HttpServer with route 
      configuration.
\item If \texttt{has\_cli}: Clap-based argument parsing dispatching 
      to subcommands.
\item If both: main.rs has CLI parsing that includes a ``serve'' 
      subcommand starting the HTTP server.
\item If \texttt{has\_library} only: no main.rs generated 
      (lib.rs with pub API).
\end{itemize}

\paragraph{\texttt{generate\_structural\_content()}} in 
\texttt{staged\_closure.rs} already dispatches by path. It now 
receives \texttt{\&self.plan.blueprint} and passes it to 
\texttt{static\_files} and \texttt{structural} generators that 
need it.

\subsubsection{Modified Files}

\begin{longtable}{lp{10cm}}
\toprule
\textbf{File} & \textbf{Change} \\
\midrule
\texttt{isls-forge-llm/src/static\_files.rs}
  & \texttt{generate\_cargo\_toml(spec, bp)}, 
    \texttt{generate\_dockerfile(spec, bp)}, 
    \texttt{generate\_docker\_compose(spec, bp)} take blueprint. \\
\texttt{isls-forge-llm/src/structural.rs}
  & \texttt{generate\_main\_rs(spec, bp)} adapts entry point.
    \texttt{generate\_for\_path(path, spec, bp)} receives blueprint.\\
\texttt{isls-forge-llm/src/staged\_closure.rs}
  & \texttt{generate\_structural\_content()} passes blueprint. \\
\bottomrule
\end{longtable}

\subsubsection{W3 Verification}

\begin{enumerate}[nosep]
\item \texttt{cargo test --workspace} --- all tests pass.
\item Unit test: \texttt{generate\_cargo\_toml()} with web blueprint 
      includes actix-web. Without: doesn't.
\item Unit test: \texttt{generate\_cargo\_toml()} with CLI blueprint 
      includes clap.
\item Unit test: \texttt{generate\_main\_rs()} with web blueprint 
      contains \texttt{HttpServer}.
\item Unit test: \texttt{generate\_main\_rs()} with CLI blueprint 
      contains \texttt{clap::Parser}.
\item Regression: run \texttt{isls forge-chat -m "Pet shop..."} with 
      default blueprint. Output identical to D5 (same files, same 
      structure, same compile result).
\end{enumerate}


% ──────────────────────────────────────────────────────────────────────────────
\subsection{W4: Pipeline Integration}
\label{sec:w4}
% ──────────────────────────────────────────────────────────────────────────────

\subsubsection{Goal}

Wire the blueprint derivation into the generation pipeline so that 
\texttt{forge-chat} and \texttt{forge-v2} automatically derive the 
blueprint from the user's description via norm matching.

\subsubsection{Flow}

\begin{lstlisting}[style=bash]
Description
  |
  v
NormRegistry::match_description()
  |
  v
activated norms (entity norms + infrastructure norms)
  |
  v
blueprint::derive_blueprint(&activated)
  |
  v
ForgePlan { spec, norm_ids, blueprint }
  |
  v
CodegenHdag::build(&spec, &blueprint)
  |
  v
S0-S7 pipeline (unchanged)
\end{lstlisting}

\subsubsection{Changes}

In \texttt{isls-chat}: after extracting entities and matching norms
(D4 enrichment), call \texttt{derive\_blueprint()} and attach the
blueprint to the \texttt{ForgePlan}.

In \texttt{isls-cli::cmd\_forge\_v2()}: when building ForgePlan from 
TOML, also match the description against infrastructure norms to 
derive a blueprint.

\subsubsection{Backward Compatibility}

If no infrastructure norms activate (old TOML files that don't have 
descriptions triggering CLI/library norms), the default web-app 
blueprint is used. All existing TOML files produce identical output.

\subsubsection{W4 Verification}

\begin{enumerate}[nosep]
\item Run all 6 D4 pipeline domains (petshop, hotel, clinic, gym, 
      school, spa) with the new pipeline. All must compile.
\item Verify generated file trees are identical to D5 output.
\end{enumerate}


% ──────────────────────────────────────────────────────────────────────────────
\subsection{W5: ISLS Studio Self-Description + forge-self}
\label{sec:w5}
% ──────────────────────────────────────────────────────────────────────────────

\subsubsection{Goal}

Create the ISLS Studio TOML and the \texttt{isls forge-self} command.
The TOML describes ISLS's own entities. The description activates
web-infrastructure norms, producing a web-app blueprint.

\subsubsection{New Files}

\begin{longtable}{lp{10cm}}
\toprule
\textbf{File} & \textbf{Purpose} \\
\midrule
\texttt{examples/isls\_studio.toml}
  & ISLS self-description: 5 entities (Norm, NormCandidate, 
    Generation, ScrapeJob, NormWiring) + User. \\
\bottomrule
\end{longtable}

The TOML uses the same format as all other ISLS requirements files. 
The \texttt{[app] description} field triggers web-infrastructure 
norms through keyword matching (``management interface'', ``dashboard'').

\texttt{isls forge-self} is 30--50 lines in \texttt{main.rs} that 
delegates to \texttt{forge-v2} with hardcoded TOML path. After 
generation, it triggers D4 observation (self-scrape).

\subsubsection{Norm Seed in Migration}

Extend \texttt{structural.rs::generate\_migration()} to detect a 
``Norm'' entity with \texttt{is\_builtin} field. If present, append 
24 INSERT statements for builtin norms. Conditional --- doesn't 
affect non-ISLS apps. Compliant with Rule~8 (seed columns from 
EntityDef).

\subsubsection{W5 Verification}

\begin{enumerate}[nosep]
\item \texttt{isls forge-self --mock-oracle} produces output with 
      correct structure.
\item \texttt{isls forge-self --api-key \$KEY} generates full Studio.
\item Generated \texttt{backend/} compiles with \texttt{cargo build}.
\item Migration contains 6 CREATE TABLE + 24 norm seed INSERTs.
\item \texttt{isls norms candidates} shows ``isls-studio'' domain.
\end{enumerate}


% ──────────────────────────────────────────────────────────────────────────────
\subsection{W6: Full Verification}
\label{sec:w6}
% ──────────────────────────────────────────────────────────────────────────────

\subsubsection{Goal}

Prove backward compatibility AND forward capability.

\subsubsection{Backward Compatibility Tests}

All 6 existing domains must still compile:

\begin{lstlisting}[style=bash]
isls forge-chat -m "Pet shop with animals, owners, appointments" \
    --api-key $KEY --output ./output/d6-test/petshop
cd output/d6-test/petshop/backend && cargo build

# Repeat for hotel, clinic, gym, school, spa
\end{lstlisting}

\subsubsection{Forward Capability: ISLS Studio}

\begin{lstlisting}[style=bash]
isls forge-self --api-key $KEY
cd output/isls-studio/backend && cargo build
# Docker optional: docker-compose up -d
\end{lstlisting}

\subsubsection{Norm System Integrity}

\begin{lstlisting}[style=bash]
isls norms list          # 24+ builtins + 7 infra + auto
isls norms stats         # Infra norms visible
isls norms candidates    # isls-studio domain observed
\end{lstlisting}


% ══════════════════════════════════════════════════════════════════════════════
\section{File Inventory}
% ══════════════════════════════════════════════════════════════════════════════

\subsection{New Files}

\begin{longtable}{llp{7cm}}
\toprule
\textbf{Crate} & \textbf{File} & \textbf{Purpose} \\
\midrule
isls-forge-llm & \texttt{src/blueprint.rs}
  & InfraBlueprint struct + derivation \\
isls & \texttt{examples/isls\_studio.toml}
  & Self-description TOML \\
\bottomrule
\end{longtable}

\subsection{Modified Files}

\begin{longtable}{llp{7cm}}
\toprule
\textbf{Crate} & \textbf{File} & \textbf{Change} \\
\midrule
isls-forge-llm & \texttt{src/lib.rs}
  & Add \texttt{pub mod blueprint}, blueprint field on ForgePlan \\
isls-forge-llm & \texttt{src/hdag.rs}
  & Refactor \texttt{build()} to take blueprint \\
isls-forge-llm & \texttt{src/static\_files.rs}
  & Blueprint-aware Cargo.toml, Dockerfile, compose \\
isls-forge-llm & \texttt{src/structural.rs}
  & Blueprint-aware main.rs, norm seed in migration \\
isls-forge-llm & \texttt{src/staged\_closure.rs}
  & Pass blueprint through pipeline \\
isls-forge-llm & \texttt{src/forge.rs}
  & Pass blueprint through LlmForge \\
isls-norms & \texttt{src/catalog.rs}
  & Add 7 infrastructure norms \\
isls-norms & \texttt{src/types.rs}
  & Add NormLevel::Atom \\
isls-cli & \texttt{src/main.rs}
  & Add forge-self, blueprint derivation in forge paths \\
isls-chat & \texttt{src/lib.rs} or \texttt{extract.rs}
  & Attach blueprint to ForgePlan after norm matching \\
\bottomrule
\end{longtable}


% ══════════════════════════════════════════════════════════════════════════════
\section{The 12 Rules --- D6 Compliance}
% ══════════════════════════════════════════════════════════════════════════════

\begin{longtable}{clp{9cm}}
\toprule
\textbf{\#} & \textbf{Status} & \textbf{D6 Impact} \\
\midrule
1  & \color{isls-green}\checkmark & LLM surface unchanged (query files 
     only). Blueprint controls which structural files exist, not which 
     are LLM-generated. \\
2  & \color{isls-green}\checkmark & mod.rs files generated for 
     whichever directories the blueprint creates. \\
3  & \color{isls-green}\checkmark & Function naming convention 
     unchanged. \\
4  & \color{isls-green}\checkmark & AuthUser still parameter (when 
     auth is active). \\
5  & \color{isls-green}\checkmark & Owned payloads. \\
6  & \color{isls-green}\checkmark & No actix\_web in queries. \\
7  & \color{isls-green}\checkmark & bcrypt/tracing when relevant. \\
8  & \color{isls-green}\checkmark & Seed INSERT uses EntityDef 
     columns. \\
9  & \color{isls-green}\checkmark & \texttt{rust:1.88-slim} when Docker 
     active. \\
10 & \color{isls-green}\checkmark & Blueprint derivation failure $\to$ 
     default web-app. \\
11 & \color{isls-green}\checkmark & ProvidedSymbol content matches 
     structural output. \\
12 & \color{isls-green}\checkmark & LLM prompts still include full 
     type context. \\
\bottomrule
\end{longtable}


% ══════════════════════════════════════════════════════════════════════════════
\section{Why This Is Not Categorization}
% ══════════════════════════════════════════════════════════════════════════════

The infrastructure norms are not categories. They are \emph{capabilities}
that compose freely:

\begin{itemize}[nosep]
\item ``Restaurant app'' $\to$ WEB + DB + AUTH + FRONTEND + DOCKER 
      = Full-stack web app with Docker deployment.
\item ``CLI tool for CSV conversion'' $\to$ CLI = Binary with clap 
      parsing, no server, no database.
\item ``Data processing service with REST API'' $\to$ WEB + DB + CLI 
      = Server with CLI management, no frontend.
\item ``Math library'' $\to$ LIB = Pure library crate, no binary, 
      no server.
\item ``Desktop app with local database'' $\to$ CLI + DB(sqlite) 
      = Binary with embedded database, no server.
\end{itemize}

No combination is predefined. The norms activate from the description.
The blueprint composes from the norms. The HDAG builds from the 
blueprint. Each layer derives from the previous --- nothing is
hardcoded.

If tomorrow a new infrastructure norm appears (e.g. 
ISLS-NORM-INFRA-GRPC for gRPC servers), it adds a 
\texttt{has\_grpc} flag to the blueprint, a conditional block in 
\texttt{build()}, and a dependency generator. No existing code changes.
The system grows by accretion, not by modification.


% ══════════════════════════════════════════════════════════════════════════════
\section{Constraints for Claude Code}
% ══════════════════════════════════════════════════════════════════════════════

\begin{enumerate}
\item \textbf{Backward compatibility is non-negotiable.} All D1--D5 
      apps must generate identically after D6. Test this with the 
      pipeline bat.
\item \textbf{Keep \texttt{build\_legacy()}} as a comparison 
      function. A test MUST verify that 
      \texttt{build(spec, default\_web\_bp())} produces equivalent 
      output to \texttt{build\_legacy(spec)}.
\item \textbf{InfraBlueprint is pure data.} No methods that generate 
      code. It describes, it doesn't do.
\item \textbf{Infrastructure norms have empty NormLayers.} They 
      influence the blueprint, not entity-level artifacts.
\item \textbf{Default fallback is web-app.} If no infrastructure 
      norms activate, the system behaves exactly as D5.
\item \textbf{Do not modify} isls-norms/src/types.rs beyond adding 
      \texttt{NormLevel::Atom}.
\item \textbf{Do not modify} isls-reader, isls-code-topo, 
      isls-hypercube, isls-types, isls-agent.
\item \textbf{Do not modify} the observation pipeline (D4) or scraping 
      pipeline (D5).
\item All new code MUST compile with \texttt{cargo build --workspace}.
\item \texttt{cargo test --workspace} --- all existing tests pass 
      (181+), plus new blueprint tests.
\item Branch: \texttt{claude/implement-isls-d6-spec}.
\item Commit format: \texttt{D6/W\{n\}: <description>}.
\end{enumerate}


% ══════════════════════════════════════════════════════════════════════════════
\section{Dependency Graph}
% ══════════════════════════════════════════════════════════════════════════════

\begin{verbatim}
W1 (InfraBlueprint + Infrastructure Norms)
  |
  v
W2 (Refactor CodegenHdag::build)
  |
  v
W3 (Adaptive Structural Generators)
  |
  v
W4 (Pipeline Integration)
  |
  v
W5 (ISLS Studio TOML + forge-self + Norm Seed)
  |
  v
W6 (Full Verification)
\end{verbatim}

Strictly sequential. Each WP depends on all previous.
Recommended implementation order: W1 $\to$ W2 $\to$ W3 $\to$ W4 
$\to$ W5 $\to$ W6.


% ══════════════════════════════════════════════════════════════════════════════
\section{Acceptance Criteria}
% ══════════════════════════════════════════════════════════════════════════════

\begin{longtable}{clp{9cm}}
\toprule
\textbf{\#} & \textbf{Pass} & \textbf{Criterion} \\
\midrule
1 & & \texttt{cargo build --workspace} compiles. \\
2 & & \texttt{cargo test --workspace} passes (181+ existing + new). \\
3 & & Comparison test: \texttt{build(spec, web\_bp) == build\_legacy(spec)}.\\
4 & & All 6 D4 pipeline apps compile with new pipeline. \\
5 & & \texttt{isls forge-self --mock-oracle} produces output. \\
6 & & \texttt{isls forge-self --api-key \$KEY} generates Studio. \\
7 & & Generated Studio compiles with \texttt{cargo build}. \\
8 & & Migration contains 6 CREATE TABLE + 24 norm seed. \\
9 & & \texttt{isls norms list} shows 7 infrastructure norms. \\
10 & & \texttt{isls norms candidates} shows isls-studio observation. \\
\bottomrule
\end{longtable}

D6 is passed when criteria 1--4 and 6--7 are fulfilled (backward 
compat proven, Studio generated and compiled). Criteria 5, 8--10 
are secondary.


% ══════════════════════════════════════════════════════════════════════════════
\section{Verification Sequence (Complete)}
% ══════════════════════════════════════════════════════════════════════════════

\begin{enumerate}
\item \texttt{cargo build --workspace}
\item \texttt{cargo test --workspace}
\item Run D4 pipeline bat --- all 6 apps compile.
\item \texttt{isls norms list} --- 24 + 7 infra norms visible.
\item \texttt{isls forge-self --mock-oracle} --- output created.
\item \texttt{isls forge-self --api-key \$KEY} --- Studio generated.
\item \texttt{cd output/isls-studio/backend \&\& cargo build} --- 
      compiles.
\item \texttt{isls norms candidates} --- isls-studio observed.
\end{enumerate}

\vfill
\begin{center}
\color{isls-gray}\rule{0.5\textwidth}{0.4pt}\\[0.5em]
{\small End of D6 Specification.\\[0.3em]
Die Infrastruktur ist keine Kategorie.\\
Sie emergiert aus den Normen.\\
Der Spielautomat hat seine Walzen.}
\end{center}

\end{document}
